\clearpage
\fancyhead[L]{\small SONNENFINSTERNISS. PSYCHE.}
\SourceLeaf{484}
\begin{center}
\rule{0.55\linewidth}{0.4pt}\\[2pt]
Astronomische Nachrichten. Band XXXIII. Nr. 771. Seite 47. 1851 August 28.\\[2pt]
\rule{0.55\linewidth}{0.4pt}
\end{center}
\begin{center}Beobachtung der \emph{Sonnenfinsterniss} 1851 Juli 28. auf der Göttinger Sternwarte.\end{center}

\begin{center}
\begin{tabular}{lrl}
Anfang & $2^{\rm u}\ 53^{\rm m}\ 37^{\rm s}4$ & \textsc{Gauss} \\
 & \mbox{---\quad---}\quad $39.4$ & \textsc{Westphal} \\
 & \mbox{---\quad---}\quad $40.4$ & \textsc{Klinkerfues} \\
\end{tabular}
\end{center}
\begin{center}
\begin{tabular}{lrl}
Ende & $5^{\rm u}\ 0^{\rm m}\ 0^{\rm s}6$ & \textsc{Gauss} \\
 & \mbox{---\quad---}\quad $1.6$ & \textsc{Westphal} \\
 & \mbox{---\quad---}\quad $2.6$ & \textsc{Klinkerfues} \\
\end{tabular}
\end{center}

\clearpage
\hypertarget{native-484-psyche}{}
\begin{center}
\rule{0.55\linewidth}{0.4pt}\\[2pt]
Astronomische Nachrichten. Band XXXV. Nr. 818. Seite 17. 1852 Juli 27.\\[2pt]
\rule{0.55\linewidth}{0.4pt}
\end{center}
\begin{center}Elemente der \emph{Psyche} nebst Bemerkungen über die Bestimmung der Bahn dieses Planeten.\end{center}

Für die Psyche ist, wegen der so sehr geringen Neigung der Bahn, eine Bestimmung aus drei Beobachtungen durchaus nicht geeignet scharfe Resultate zu geben, wie dies auch die von verschiedenen Rechnern auf diese Art erhaltenen Elemente bestätigen. Ich habe daher Herrn \textsc{Klinkerfues} veranlasst, eine neue Bestimmung auf vier Oerter zu gründen, nach der in der Th. M. C. C. entwickelten, aber wie es scheint bisher wenig oder gar nicht zur Anwendung gebrachten Methode. Ich lege seine Elemente hier bei, zugleich mit der Vergleichung mit allen uns bekannt gewordenen Beobachtungen.

Interessant ist die nahe Uebereinstimmung mit \textsc{Hind}'s Schätzung vom 29. Januar, so dass an der Identität des damals von \textsc{Hind} gesehenen Sterns mit dem 7 Wochen später von \textsc{Gasparis} entdeckten Planeten nicht mehr zu zweifeln ist. Wenn Herr \textsc{Hind} etwas näheres über seine damalige Aufzeichnung mittheilen wollte, namentlich über die Stunde und über den Grad von Genauigkeit, welchen er seiner Schätzung beilegt, so würde wahrscheinlich die Mitzuziehung dieser Wahrnehmung zu einer neuen Berechnung die Genauigkeit des Resultats noch merklich erhöhen können.

\Needspace{17\baselineskip}
\begin{center}
Elemente der Psyche.\\
Epoche: März 31.0 mittlere Zeit Berlin.
\end{center}
\[
\text{Mittlere Länge}=151^\circ\ 32'\ 5''56
\]
\[
\left.\begin{array}{rcl}
\pi &=& 6\ 20\ 2.86 \\
\AscendingNode &=& 150\ 33\ 5.86
\end{array}\right\}\quad\text{Mittl. Aequin. 1852.0}
\]
\[
\begin{array}{rcl}
i &=& 3\ 2\ 16.06 \\
\varphi &=& 6\ 38\ 32.80 \\
\log e &=& 9.0632320 \\
\log a &=& 0.4693195 \\
\mu &=& 701''4993
\end{array}
\]

Vorstehende Elemente sind aus 4 Normal-Oertern von März 19, April 14, April 30 und Mai 14 berechnet.

Unter den unten [hier nicht abgedruckten] mit diesem Systeme verglichenen Beobachtungen befindet sich auch die von Herrn \textsc{Hind} am 29. Januar gemachte Schätzung. Bei Unkenntniss der Zeit, welcher sie entspricht, ist angenommen, dieselbe sei etwa Berliner Mitternacht gewesen. Für diese Zeit ist die stündliche Bewegung des Planeten $22''75$ nach West und $10''25$ nach Nord.

\clearpage
\fancyhead[L]{\small ASTRONOMISCHE NACHRICHTEN.}
\SourceLeaf{485}
\noindent[Die in den Astronomischen Nachrichten abgedruckten Auszüge aus späteren Briefen von \textsc{Gauss}

\begin{center}
\setlength{\tabcolsep}{3pt}
\begin{tabular}{llrlrlll}
nemlich & Band & 36 & Nr. & 864 & Göttingen & 1853 & Juni 12 \\
 & & 37 & & 875 & & 1853 & September 3 \\
 & & 37 & & 885 & & 1853 & December 2 \\
 & & 37 & & 888 & & 1853 & December 16 \\
 & & 38 & & 910 & & 1854 & Juni 30 \\
\end{tabular}
\end{center}

\noindent enthalten Mittheilungen der von \textsc{Klinkerfues}, \textsc{Westphal} und \textsc{Winnecke} auf der Sternwarte in Göttingen zu jener Zeit ausgeführten Beobachtungen so wie der Entdeckungen von Cometen durch \textsc{Klinkerfues}.

\begin{flushright}\textsc{Schering}.]\end{flushright}
